%==============================================================================
% Section 2: Related Work
% Verbatim from Nckh-official_2.docx (P19-P28).
% The source uses IEEE-style letter labels ("B.", "C.", "D.") on the subsection
% titles; they are dropped here because llncs.cls numbers subsections
% automatically (2.1, 2.2, ...), which reproduces the same numbering.
%==============================================================================
\section{Related Work}
\label{sec:related}

% ---- 2.1 --------------------------------------------------------------------
\subsection{Deep Sequence Modeling: From RNNs to Mamba}
\label{subsec:related-sequence}

In Section~\ref{sec:introduction}, we present an overview of the general path
from recurrent architectures to Transformers and selective state space models.
This evolution has resulted in several Mamba-based implementations
(MambaLithium, SambaMixer and iMamba variants) in battery prognostics claiming
better accuracy and efficiency for unified SOH, SOC and RUL estimation than
recurrent and attention-based baselines~\cite{huang2025,olalde2024}.

% ---- 2.2 --------------------------------------------------------------------
\subsection{Health Indicators and Multimodal Frequency Conditioning}
\label{subsec:related-indicators}

These can be telemetry values such as raw voltage/current/temperature or more
degradation-sensitive values such as incremental capacity and differential
voltage analysis~\cite{su2024}, electrochemical impedance
spectroscopy~\cite{alsmadi2025}, and internal gas pressure~\cite{wang2025}
enabling multi-feature approaches instead of only time-domain
ones~\cite{yang2023}. It is difficult to fuse such multi-domain information with
high-dimensional temporal sequences because of mismatched statistical
distributions and manifolds. Modal conflicts and suboptimal gradient flow become
common consequences of naive concatenation~\cite{meng2026}. Architectures like
Hybrid Sensing Synergy Architecture (HSSA) use Q-Former and cross-attention
modules to align EIS features to the temporal representation
space~\cite{meng2026}, and Feature-wise Linear Modulation (FiLM) provides a
lighter approach to adapt sequence representations directly from spectral
conditioning without specialised EIS hardware~\cite{huang2025}.

% ---- 2.3 --------------------------------------------------------------------
\subsection{Robust Diagnostics and Anomaly Detection}
\label{subsec:related-anomaly}

In order to design a robust and safe AI powered Battery Management
System~\cite{arbaoui2024}, it is needed to estimate the point-wise
state-of-health (SOH), quantify uncertainty, and detect anomalies using
state-of-the-art models. Anomalies can be caused by sensor noise, distributional
shifts or new internal faults and result in accelerated degradation. They are
detected by anomaly detection~\cite{arbaoui2024}.
% NOTE: the source manuscript cites "[34]" at this exact point, but its
% reference list has only 23 entries, so [34] does not resolve. Mapped to
% Arbaoui et al. (2024) -- the closest source by topic. Please confirm the
% intended reference and change the \cite key above if it is a different one.
Hybrid methods have been successfully applied to practical solar-grid and EV
cases such as statistical methods such as local outlier factor (LOF) and
Mahalanobis distance, and advanced frameworks such as Markov-constrained
Isolation Forest~\cite{aljohani2026}. The deployment of these powerful
diagnostic tools in efficient Mamba based
architectures~\cite{meng2026,olalde2024} marks a milestone in the development of
self-aware, responsible and deployable battery health monitoring systems.

% ---- 2.4 --------------------------------------------------------------------
\subsection{Literature Comparison Matrix}
\label{subsec:related-matrix}

% Table 1 of the source manuscript. Caption above the table and a ruled (grid)
% layout with \hline, following the Springer LNCS samplepaper.
\begin{table}[H]
  \caption{Comprehensive Literature Comparison Matrix and Structural
           Comparison}
  \label{tab:literature}
  \centering
  \scriptsize
  \linespread{0.92}\selectfont
  \tabtight
  % Note: plain spaces (not ~) before \cite inside the cells, so that a long
  % study name and its citation label can break across lines.
  \begin{tabularx}{\textwidth}{|M{1.95cm}|M{1.5cm}|Y|Y|Y|}
    \hline
    \textbf{Study} & \textbf{Method Family} & \textbf{Key Strength} &
    \textbf{Key Limitation} & \textbf{Relation to This Work} \\
    \hline
    LSTM / CNN-LSTM \cite{arbaoui2024,zhao2024}
      & Sequential / Hybrid CNN-RNN
      & Effective temporal modeling on short windows; simple implementation
      & Recurrence bottleneck; struggles with long-range dependencies and
        cross-battery generalization
      & Served as primary baseline for comparison; outperformed by proposed
        Mamba architecture \\
    \hline
    Transformer / PatchTST \cite{patel2025}
      & Attention-based Transformer
      & Strong global context capture via self-attention
      & Quadratic complexity $O(L^{2})$; high memory/latency on long sequences
        ($L=4096$); edge deployment challenges
      & Advanced baseline demonstrating limitations addressed by linear Mamba
        backbone \\
    \hline
    Unimodal Mamba (SambaMixer,
        MambaLithium) \cite{olalde2024}
      & Selective State Space Model (SSM)
      & Linear complexity $O(L)$; efficient long-sequence modeling
      & Unimodal (raw telemetry only); limited physics-informed conditioning
        and interpretability
      & Core backbone adopted and extended; this work adds spectral
        conditioning via FiLM \\
    \hline
    HSSA (Mamba $+$ EIS Multimodal) \cite{meng2026}
      & Mamba + Q-Former Multimodal Fusion
      & Multimodal fusion of discharge curves and impedance spectra; strong
        performance on NASA
      & Requires EIS data (not always available in field); higher complexity
        for edge devices
      & Closely related multimodal inspiration; this work uses spectral
        features (FFT/statistical) as more lightweight alternative
        conditioning signal \\
    \hline
  \end{tabularx}
  \srcnote{Source: Summarize of Authors, 2026}
\end{table}
